% ---- pg-224 （书页 212） ----
试问，这两家实验室的化学分析结果是否一致？若不一致，原因是什么？

检验 $H_0: \boldsymbol{\delta}'=[\delta_1,\delta_2]=[0,0]$ 的 $T^2$ 统计量可用下列成对观测值之差来构造：
\[
\begin{array}{l r r r r r r r r r r r}
d_{j1}=x_{1j1}-x_{2j1} & -19 & -22 & -18 & -27 & -4 & -10 & -14 & 17 & 9 & 4 & -19\\
d_{j2}=x_{1j2}-x_{2j2} & 12 & 10 & 42 & 15 & -1 & 11 & -4 & 60 & -2 & 10 & -7
\end{array}
\]
由此得
\[
\bar{\mathbf{d}}=\begin{bmatrix}\bar{d}_1\\\bar{d}_2\end{bmatrix}=\begin{bmatrix}-9.36\\13.27\end{bmatrix},\qquad
\mathbf{S}_d=\begin{bmatrix}199.26&88.38\\88.38&418.61\end{bmatrix}
\]
因而
\[
T^2=11[-9.36,13.27]\begin{bmatrix}0.0055&-0.0012\\-0.0012&0.0026\end{bmatrix}\begin{bmatrix}-9.36\\13.27\end{bmatrix}=13.6
\]
取 $\alpha=0.05$，求得 $[(n-1)p/(n-p)]F_{p,n-p}(\alpha)=(2\times 10/9)F_{2,9}(0.05)=9.47$，因为 $T^2=13.6>9.47$，我们拒绝 $H_0$，并得出这两家实验室的检验结果有显著差异的结论. 从数据的检查中也发现，商业实验室提供的生化氧(BOD)测量值要比州卫生实验室的测量值低，而悬浮固体(SS)含量的测量值则要比后者的测量值高. 从\hyperref[eq:6-10]{式(6-10)}可算出均值差 $\delta_1$ 和 $\delta_2$ 的 95\% 联合置信区间：

\[
\delta_1: \bar{d}_1\pm\sqrt{\frac{(n-1)p}{(n-p)}F_{p,n-p}(\alpha)}\sqrt{\frac{s_1^2}{n}}
=-9.36\pm\sqrt{9.47}\times\sqrt{\frac{199.26}{11}}\quad 即\ (-22.46,3.74)
\]

\[
\delta_2: 13.27\pm\sqrt{9.47}\sqrt{\frac{418.61}{11}}\quad 即\ (-5.71,32.25)
\]

虽然这两个 95\% 的联合置信区间均包含 0，但在 0.05 显著性水平下，还是拒绝假设 $H_0: \boldsymbol{\delta}=\mathbf{0}$. 我们能得出什么结论？

这一证据暗示的确存在差别. 点 $\boldsymbol{\delta}=\mathbf{0}$ 落在 $\boldsymbol{\delta}$ 的 95\% 置信区域之外（见练习 6.1），因而这一结果与 $T^2$ 检验的结果一致. 95\% 的联合置信系数适用于能用 $\delta_1,\delta_2$ 的任一可能的线性组合 $a_1\delta_1+a_2\delta_2$ 构造的区间集合. 与 $(a_1=1,a_2=0)$ 和 $(a_1=0,a_2=1)$ 这两种选择对应的特定区间均包含 0，但对 $a_1,a_2$ 的其他选择将产生出不包含 0 的联合置信区间（若假设 $H_0: \boldsymbol{\delta}=\mathbf{0}$ 不被拒绝，则所有的联合置信区间均包含 0）.

庞弗罗尼联合置信区间也包含 0（见练习 6.2）.

我们的分析是在 $\mathbf{D}$ 的正态性假定下进行的. 实际上，当存在一个或两个离群值时，情况将更加复杂（见练习 6.3）. 虽然这些数据可变换成更接近正态性的数据，但对这种小样本来说，要消除离群值效应实在很困难（见练习 6.4）.

本例的数值结果表明，在我们进行统计推断时，可能发生某种不正常情况. \hfill$\blacksquare$

例 6.1 中的实验人员实际上要先摇晃样品，再迅速将它分别倒入两个瓶子里以供化学分析. 这样做是很明智的，因为若是简单地把瓶颈中的样品的一半倒入一个瓶，剩余一半倒入另一个瓶，则由于沉淀之故，第二个瓶中会有较多的悬浮固体. 这样，两个实验室所处理的实验单元就不会相同，因而它们得出的结论也就不可能反映实验室的水平、测量技术等方面的差异.

每当研究人员能控制将处理分配给各实验单元的过程时，实验单元的适当配对和处理的随机分配就能提高统计推断的精度. 一些被设想为相同的实验单元之间若有差异，必须将这
